% Inclusion fragment. Requires obstruction.tex and finite_profiles.tex.
% Requires theorem, lemma, corollary, definition, remark environments.
% Bibliography key: avigad-cut.
% Scope: ordinary intuitionistic theories with existential-free axioms and
% the standard B-algebra as a model. Not arbitrary axiomatizations in B syntax.

\section{The obstruction to relational bi-interpretability}
\label{sec:relational-obstruction}

The preceding semantic obstructions concerned operation symbols realized
by terms. The finite-profile theorem also has a proof-theoretic
consequence. For the ordinary tied-seed arithmetic with literal-equation
induction, it rules out \emph{every finite-dimensional
bi-interpretation} with $\mathsf{iPRA}$, including interpretations with
arbitrary formula-defined domains, operation graphs, and quotient
equality. The additional ingredient is finite witness extraction from
intuitionistic proofs with existential-free axioms.

This conclusion concerns a specified theory, not every possible theory
in the same numerical language. Adding positive existence axioms can
invalidate the witness-extraction hypothesis. It also does not claim
that the tied-seed theory has no one-way interpretation of arithmetic
on a complicated domain. The source-side round trip required for
bi-interpretability is essential to the final argument.

\subsection{The ordinary tied-seed theory has an existential-free presentation}

Let $T_B$ use ordinary intuitionistic first-order logic with equality and
the staged $0,B$ numerical signature. It has the two computation schemas
\eqref{obs:B0}--\eqref{obs:BS}, numerical coverage
\[
              \forall x\,(x=0\vee\exists y\,x=Sy),
\]
separation $(0=S0)\to\bot$, and literal-equation induction. The latter
may be given as the universal closures of
\[
 \left(E(0,\vec p)\wedge
       \forall n\,(E(n,\vec p)\to E(Sn,\vec p))\right)
       \longrightarrow\forall n\,E(n,\vec p),
 \tag{EI axiom}\label{rel:EI-axiom}
\]
where $E$ is a literal equation. This axiom presentation and the usual
contextual induction rule are interderivable by implication and
universal introduction and elimination, with the ordinary eigenvariable
conditions.

Every axiom in this list except the displayed coverage sentence is
existential-free. The predecessor abbreviation $P$ already satisfies
\[
                         P0=0,\qquad P(Sy)=y
\]
by computation. Consequently coverage is equivalent to
\[
                      \forall x\,(x=0\vee x=S(Px)).
                      \tag{Coverage$'$}\label{rel:coverage-prime}
\]
Indeed, a witness $y$ to $x=Sy$ is equal to $Px$, and the converse uses
the explicit witness $Px$. Neither direction needs induction. Thus
$T_B$ has an equivalent existential-free axiomatization in the same
numerical language.

For the proof-theoretic argument, equality can be treated as an ordinary
binary predicate with universal axioms for reflexivity, symmetry,
transitivity, and congruence of every function symbol and atomic predicate
symbol. In the numerical $B$ language equality is the only predicate, so
its predicate-congruence instances already follow from symmetry and
transitivity. All these axioms
are existential-free; transport through arbitrary formulas is derived
by induction on formula syntax. This avoids invoking a primitive
equality rule whose motive could obscure the subformula argument.
The recursion notation is expanded into its staged first-order function
symbols before applying first-order cut elimination.

The argument below applies more generally to any ordinary intuitionistic
theory $T$ in this numerical language that has an existential-free
axiomatization and has the standard $B$-algebra as a model. An infinite
axiom scheme causes no difficulty: a finite proof uses only finitely
many axioms. Extra existential-free axioms that are true in the standard
model do not evade the result.

\subsection{Finite witnesses with an arbitrary matrix}

\begin{lemma}[Finite witness extraction]
\label{rel:finite-witnesses}
Let the language contain a closed term $0$, and let $\Gamma$ be a finite
set of existential-free formulas. If intuitionistic first-order logic
proves
\[
                      \Gamma\vdash\exists z_1\cdots\exists z_d\,
                                      R(\vec x,\vec z),
\]
then there are finitely many tuples of language terms
$\vec t_1(\vec x),\ldots,\vec t_M(\vec x)$ for which
\[
                   \Gamma\vdash
                   \bigvee_{\ell=1}^{M}R(\vec x,\vec t_\ell(\vec x)).
                   \tag{Witness list}\label{rel:witness-list}
\]
Any additional fixed parameters are retained in the terms. The formula
$R$ may contain arbitrary quantifiers, including existential quantifiers.
No uniqueness of witnesses is assumed.
\end{lemma}

\begin{proof}
Use a cut-free intuitionistic sequent proof. The required cut-elimination
theorem is standard; for a published proof see
Avigad~\cite{avigad-cut}, Theorem~7.6. Equality is handled by the
existential-free axioms just specified, so no special equality extension
of that theorem is needed.

Starting at the end sequent, trace only the branches carrying the
displayed leading existential quantifiers. Keep the finite prefix of
witness terms already selected. At a left implication rule retain its
first premise proof unchanged and follow its second premise, whose
succedent remains the tracked existential formula. At a left
disjunction rule follow both premises. The other applicable left and
structural rules preserve the tracked succedent in the usual way.

Every antecedent formula reached on these tracked branches is a
subformula instance of a formula in $\Gamma$, and is therefore
existential-free. In particular, no left existential rule can occur
there. While at least one of the displayed leading existential
quantifiers remains, an initial identity sequent cannot close the
branch: its existential succedent is not an antecedent formula. The
only applicable right introduction is existential introduction. Record
its witness term and continue with the shorter leading prefix. When
all $d$ witnesses have been introduced, stop tracing that branch and
retain the remaining proof of its matrix instance unchanged. The
matrix can be arbitrarily complicated because it is not analyzed.

If a tracked branch instead closes by falsity elimination, fill any
unselected witnesses with $0$; its resulting matrix instance follows
from the same falsity. There are finitely many branches, so this process
records finitely many complete tuples $\vec t_\ell$.

Now replace each stopped branch by its retained proof of
$R(\vec x,\vec t_\ell)$ followed by the appropriate introductions into
the common finite disjunction in~\eqref{rel:witness-list}. Reassemble
the tracked left and structural inferences. At an existential
introduction that supplied a recorded component, the finite disjunction
is already the succedent of the transformed premise, so that
introduction is omitted. At left implication inferences the unchanged
first premise and the transformed second premise reconstruct the
inference. This gives the claimed proof.

There are no newly bound numerical parameters on the tracked branches:
left existential introduction is absent, and right universal
introduction cannot have the tracked existential succedent. Any
extraneous free parameters of the original finite proof can be replaced
uniformly by $0$, after renaming local eigenvariables. Thus the witness
terms need only the retained parameters and $\vec x$.
\end{proof}

This lemma is an application of cut elimination, not a claim that
Avigad's paper contains the present arithmetic obstruction. It also
does not assert a single uniform term witness. A finite disjunction is
exactly what is needed below.

\subsection{An arity obstruction for provably total coding relations}

\begin{theorem}[No provable compression from $k$ scalars to $d$ coordinates]
\label{rel:no-compression}
Let $T$ have an existential-free axiomatization and the standard
$B$-algebra as a model. Fix positive integers $k>d$ and any formula
$R(\vec x,\vec z)$, where $\vec x$ has $k$ scalar coordinates and
$\vec z$ has $d$. Suppose $T$ proves
\[
                        \forall\vec x\,\exists\vec z\,R(\vec x,\vec z).
                        \tag{Total}\label{rel:total}
\]
Then it is impossible that, in the standard model, each output tuple
$\vec z$ occurs with at most one input tuple $\vec x$. In particular,
$T$ cannot also prove
\[
 R(\vec x,\vec z)\wedge R(\vec x\,',\vec z)
             \longrightarrow\bigwedge_{i=1}^{k}x_i=x_i'.
 \tag{Disjoint fibers}\label{rel:disjoint-fibers}
\]
The relation may have many outputs for each input.
\end{theorem}

\begin{proof}
Instantiate~\eqref{rel:total} at free variables $\vec x$ and collect
the finitely many axioms used by the proof in an existential-free
antecedent. Lemma~\ref{rel:finite-witnesses} yields a finite list of
$d$-tuples of $B$-terms $\vec t_1,\ldots,\vec t_M$ such that, for every
input in the standard model, at least one listed tuple is an $R$-output.

For each scalar coordinate term $t_{\ell j}$, apply
Theorem~\ref{fc:B-theorem}. Let $q_{\ell j}$ be the finite number of
its possible unary profiles. On the grid
\[
                         X_N=\{0,1,\ldots,N-1\}^{k},
\]
that term has at most $q_{\ell j}N$ possible values: its active argument
has at most $N$ values, and there are at most $q_{\ell j}$ profiles.
Constant terms can be treated by adding an unused active argument.
Thus the vector $\vec t_\ell$ has at most
\[
                         N^d\prod_{j=1}^{d}q_{\ell j}
\]
possible values on $X_N$. The union of the images of all listed tuples
has size at most $CN^d$, where
\[
                         C=\sum_{\ell=1}^{M}\prod_{j=1}^{d}q_{\ell j}
\]
is fixed independently of $N$.

If each output belonged to at most one input, the $N^k$ inputs in $X_N$
would require at least $N^k$ distinct outputs in this union. This is
impossible for a sufficiently large standard $N$, since $k>d$.
Taking $N>C$ already suffices. Soundness in the standard model turns any
proved instance of~\eqref{rel:disjoint-fibers} into exactly the forbidden
semantic condition.
\end{proof}

\begin{corollary}[Scalar pairing totality cannot be proved]
\label{rel:no-scalar-pairing}
In $T_B$, no formula can be proved to code every pair of source numbers
by a single source number with unique decoding. This includes any
proposed pairing relation
\[
                         \pi_1(c)=a\ \wedge\ \pi_2(c)=b
\]
using fixed unary $B$-terms as projections. The conclusion remains true
even if the projection relations are arbitrary formulas, provided each
code determines at most one pair.
\end{corollary}

Thus even a future proof of unary primitive-recursive completeness of
the $B$ term language would not establish the required pairing totality
inside $T_B$. The obstacle is proof-theoretic as well as term-theoretic.

\subsection{Why bi-interpretability supplies forbidden compression}

\begin{theorem}[No finite-dimensional bi-interpretation with $\mathsf{iPRA}$]
\label{rel:no-biinterpretation}
Let $T$ have an existential-free axiomatization in the $B$ numerical
language and have the standard $B$-algebra as a model. Then $T$ and full
first-order $\mathsf{iPRA}$ are not bi-interpretable by ordinary
finite-dimensional first-order interpretations. The interpreted domains,
operation graphs, and quotient equality formulas may be arbitrary.

The statement is parameter-free, and also applies when fixed parameters
are retained uniformly in the interpretation proofs without adding
non-existential-free assumptions to $T$.
\end{theorem}

\begin{proof}
Suppose that $I$ interprets arithmetic $A=\mathsf{iPRA}$ in $T$, and
$J$ interprets $T$ in $A$. Write
\[
                          I:A\rightsquigarrow T,
                          \qquad J:T\rightsquigarrow A.
\]
Let $I$ have dimension $d$, with numerical domain
$D\subseteq\mathbb N^d$ and interpreted arithmetic equality $E$.
Let $J$ have dimension $e$ over the arithmetic domain. A zero-dimensional
domain has only finitely many tuples, so it cannot participate in the
assumed source-side isomorphism with the infinite standard source
domain. We may therefore take $d,e\ge1$.

The composite $I\circ J$ interprets source numbers by suitable
$e$-tuples of interpreted arithmetic elements, represented in the
source by $de$ numerical coordinates. The source-side bi-interpretation
isomorphism provides a $T$-definable relation $\eta(a,\vec u)$, with
$\vec u\in D^e$, that is provably total for every source number $a$
and is injective modulo the composite interpreted source equality.
In particular, $T$ proves
\[
 \forall a\,\exists\vec u\,\eta(a,\vec u),
 \qquad
 \eta(a,\vec u)\wedge\eta(a',\vec u\,')\wedge
       E_{J}^{I}(\vec u,\vec u\,')\longrightarrow a=a'.
 \tag{Source round trip}\label{rel:source-round-trip}
\]
Domain conditions are included in $\eta$. Here $E_{J}^{I}$ denotes
the interpretation under $I$ of the equality relation used by $J$.
No restrictions are placed on the quantifiers in $\eta$ or the domain
formulas. Only this source-side part of bi-interpretability will be used.

Fix an external positive integer $k$. Arithmetic $A$ has
primitive-recursive injective tupling for $ke$ numerical inputs, with
primitive-recursive projections; its decoding equations are theorems
of $A$. Interpret that tupling under $I$. It gives a provably total
formula-defined coding relation from $(D/E)^{ke}$ to $D/E$, injective
on its input tuples.

Define $R_k(a_1,\ldots,a_k,\vec c)$ in $T$ as follows: there exist
$\eta$-representatives $\vec u_i\in D^e$ of all $a_i$, and the
interpreted arithmetic tupling relation codes the $ke$ arithmetic
components of $(\vec u_1,\ldots,\vec u_k)$ by $\vec c\in D$.
This is one first-order formula. Its matrix may contain arbitrary
quantifiers. From the first clause of~\eqref{rel:source-round-trip} and
the totality of arithmetic tupling, ordinary finite existential
reasoning proves
\[
                      T\vdash\forall a_1\cdots\forall a_k\,
                         \exists c_1\cdots\exists c_d\,
                              R_k(\vec a,\vec c).
                      \tag{Compression totality}\label{rel:compression-totality}
\]
There is no interpreted-domain assumption on the input scalars $a_i$:
they range over the entire source domain. No form of choice is required;
only the finitely many witnesses in this fixed formula are introduced.

Suppose that $R_k(\vec a,\vec c)$ and
$R_k(\vec a\,',\vec c)$ hold. The same output representative
$\vec c$ represents the same arithmetic element, by reflexivity of $E$
on $D$. Injectivity of interpreted arithmetic tupling therefore gives
componentwise $E$-equality of all corresponding arithmetic components
of $\vec u_i$ and $\vec u_i'$. Interpretations respect equality, so
reflexivity and congruence for the interpreted $J$-equality imply
\[
                         E_{J}^{I}(\vec u_i,\vec u_i')
                         \qquad(1\le i\le k).
\]
The second clause of~\eqref{rel:source-round-trip} gives
$a_i=a_i'$ for every $i$. These are all legitimate interpreted
first-order consequences, so $T$ proves that the output fibers of
$R_k$ are disjoint in exactly the sense of
\eqref{rel:disjoint-fibers}. Surjectivity of $\eta$, uniqueness of the
chosen representatives, and uniqueness of the raw output tuples are
not needed.

Now choose $k=d+1$. Formula~\eqref{rel:compression-totality} is a
provably total relation from $k$ unrestricted source scalars to $d$
source coordinates with disjoint output fibers. This contradicts
Theorem~\ref{rel:no-compression}.
\end{proof}

\begin{corollary}[Provably inhabited finite parameter families]
\label{rel:parameter-families}
The exclusion also applies to a finite parameterized bi-interpretation
when $T$ proves that its parameter condition is inhabited and proves the
required interpretation and source-round-trip facts uniformly under
that condition.
\end{corollary}

\begin{proof}
Let $\vec p$ be the finite tuple of source parameters, of length $r$,
and let its condition be $\Theta(\vec p)$. Assume
$T\vdash\exists\vec p\,\Theta(\vec p)$ and that all the preceding
interpretation and isomorphism facts are proved uniformly under
$\Theta(\vec p)$. Collect any additional finite parameter choices
needed for those facts into the same tuple and condition.

The preceding construction, with these parameters retained, proves
the totality and disjoint-fiber properties of $R_{k,\vec p}$ under
$\Theta(\vec p)$. Introduce the parameters as part of the output:
\[
 \widehat R_k(\vec a,\vec p,\vec c)
       :=\Theta(\vec p)\wedge R_{k,\vec p}(\vec a,\vec c).
\]
Using the proved inhabitedness first, $T$ proves the unconditional
sentence
\[
            \forall\vec a\,\exists\vec p\,\exists\vec c\,
                  \widehat R_k(\vec a,\vec p,\vec c).
\]
An identical output pair $(\vec p,\vec c)$ uses identical parameter
values. The disjoint-fiber argument for $R_{k,\vec p}$ therefore still
gives equality of the input tuples. The number of output coordinates is
now $d+r$, so choose $k=d+r+1$ and apply
Theorem~\ref{rel:no-compression}. No restriction on the quantifiers in
$\Theta$ is needed: it is inside the output matrix, rather than an
extra assumption of the existential-free axiom theory.
\end{proof}

An externally chosen, model-specific parameter condition that $T$ does
not prove to be inhabited is not covered by this extension. Fixed
parameters for which the interpretation facts are proved without extra
assumptions are already covered by Theorem~\ref{rel:no-biinterpretation}.

The proof explains why arbitrary positive existential domain formulas
do not rescue \emph{bi}-interpretability: they occur inside the matrix
$R_k$, while the source-side isomorphism removes every domain condition
from its unrestricted source inputs. Finite witness extraction makes no
syntactic restriction on that matrix.

The theorem is a completed negative result for the ordinary tied-seed
candidate and for the stated class of its extensions. It does not
provide a smaller positive axiomatization of $\mathsf{iPRA}$, does not
rule out equiconsistency alone, and does not claim that every possible
axiomatization in the $B$ language admits existential-free axioms.
It also does not automatically extend to additional bounded numerical
primitives: they preserve the weaker coordinate bound, but need not
preserve the finite-profile invariant used here.
